% moved/renamed from 'Difference Equations Summary.tex'
\documentclass[a4paper,12pt]{article}

\usepackage{fontspec}
\usepackage{xeCJK}
\usepackage{amsmath,amssymb,geometry,graphicx}
\geometry{a4paper, margin=1in}
\setlength{\parskip}{0.8em}
\setlength{\parindent}{0em}
\setmainfont{Times New Roman}
\setCJKmainfont{標楷體}
% 若系統無此字型，可改用：
% \setCJKmainfont{DFKai-SB} % Windows 備選
% \setCJKmainfont{Noto Serif CJK TC} % 跨平台常見
\XeTeXlinebreakskip = 0pt plus 1pt


\title{Difference Equations Summary}
\author{李辰安}
\date{\today}


\begin{document}
\maketitle
\section*{1. 問題背景}
考慮簡諧運動的微分方程：
\[
y'' = -y,
\]
其解為 $y(t) = \cos t$，在 $(y, y')$ 平面上是一個圓。若要在電腦上模擬此運動，必須將連續時間轉為離散時間 $t_n = n\Delta t$，以差分近似取代微分：
\[
\frac{Y_{n+1} - 2Y_n + Y_{n-1}}{(\Delta t)^2} \approx -Y_n.
\]


\section*{2. 三種差分方法}
以下為將 $y'' = -y$ 直接離散化的三種方式（以 RHS 上 $Y$ 的取樣點命名；注意這些二階差分法與第三節「先改寫為一階系統再用 Euler 法」不同）：
\begin{center}
\begin{tabular}{|c|c|c|}
\hline
方法（取樣點） & 近似公式 & 數值行為（Eigenvalue 性質） \\
\hline
Forward（用 $Y_{n-1}$） & $Y_{n+1} = 2Y_n - Y_{n-1} - (\Delta t)^2 Y_{n-1}$ & $|\lambda| > 1$，螺旋發散 \\
Centered（用 $Y_{n}$；Leapfrog） & $Y_{n+1} = 2Y_n - Y_{n-1} - (\Delta t)^2 Y_n$ & 對 SHO：$|\lambda| = 1$ 需 $\Delta t \le 2$ \\
Backward（用 $Y_{n+1}$） & $Y_{n+1} = 2Y_n - Y_{n-1} - (\Delta t)^2 Y_{n+1}$ & $|\lambda| < 1$，螺旋收縮 \\
\hline
\end{tabular}
\end{center}
註：對 $y'' + y = 0$，Centered（Leapfrog）在 $\Delta t \le 2$ 時滿足 $|\lambda|=1$（界內穩定）；若 $\Delta t>2$ 則不穩定。


\section*{3. 化為一階系統}
定義 $Z_n \approx Y'(t_n)$，並設
\[
U_n = \begin{bmatrix} Y_n \\ Z_n \end{bmatrix}.
\]
對於 Forward Euler（先將 $y''=-y$ 寫為 $u' = Ju$，$J=\begin{bmatrix}0&1\\-1&0\end{bmatrix}$），有
\[
U_{n+1} = (I + \Delta t\, J)\, U_n, \quad
A := I + \Delta t\, J = \begin{bmatrix} 1 & \Delta t \\ -\Delta t & 1 \end{bmatrix}.
\]
Backward Euler 對同一一階系統的一步矩陣為
\[
U_{n+1} = (I - \Delta t\, J)^{-1} U_n \;=:\; B.
\]


\section*{4. 特徵值分析}
矩陣 $A$ 的特徵值為：
\[
\lambda(A) = 1 \pm i\Delta t, \quad |\lambda(A)| = \sqrt{1 + (\Delta t)^2} > 1.
\]
因此 $(Y_n, Z_n)$ 的軌跡會逐漸往外螺旋。對於 Backward Euler，
\[
B = (I - \Delta t\, J)^{-1} \;\;\Rightarrow\;\; \lambda(B) = \frac{1}{1 \mp i\Delta t}, \quad |\lambda(B)| = \frac{1}{\sqrt{1+(\Delta t)^2}} < 1,
\]
故軌跡往內收縮。中心差分（Leapfrog）在滿足 $\Delta t \le 2$ 的條件下有 $|\lambda|=1$，近似能量守恆，軌跡接近圓形（帶相位誤差）。


\section*{5. 幾何詮釋}
真實解 $u(t) = (\cos t, -\sin t)$ 恆在單位圓上。三種方法的幾何行為如下：
\begin{itemize}
\item \textbf{Forward Euler（前向法，一階）}: 軌跡向外膨脹，能量增加。
\item \textbf{Backward Euler（後向法，一階）}: 軌跡向內收縮，能量減少。
\item \textbf{Centered／Leapfrog（二階差分）}: 在 $\Delta t \le 2$ 時軌跡保持在圓附近，最穩定。
\end{itemize}


\section*{6. 數值穩定性意義}
此例展示了\textbf{數值穩定性（numerical stability）}的概念。系統長期行為取決於特徵值的模長：
\begin{itemize}
\item 若 $|\lambda|>1$：系統不穩定，誤差放大。
\item 若 $|\lambda|=1$：系統邊界穩定，能量（近似）守恆，但可能有相位漂移。
\item 若 $|\lambda|<1$：系統耗散，能量衰減。
\end{itemize}
因此，中心差分（Leapfrog）法常用於模擬振動、天體運動或分子動力學等守恆系統，但需選擇合適步長（如 $\Delta t \le 2$ 對 SHO）。


\section*{7. 學習重點整理}
\begin{enumerate}
\item 了解如何將 $y''=-y$ 轉為差分形式（索引一致性很重要）。
\item 學會將二階方程改寫為一階系統 $U_{n+1}=AU_n$ 或 $U_{n+1}=BU_n$。
\item 透過特徵值分析長期穩定性（\(|\lambda|\) 大小）。
\item 建立幾何直覺：模長決定螺旋方向與能量變化。
\item 認識 Leapfrog 法的時間對稱性、穩定域與近似能量守恆特性。
\end{enumerate}

\end{document}

